% Part: sets-functions-relations
% Chapter: size-of-sets
% Section: non-enumerability
% Hindi translation (hi-Deva-IN) of Open Logic commit 9620cc73f9c8e0ad003c514a5d3748f29611c4c0.

\documentclass[../../../include/open-logic-section]{subfiles}

\begin{document}

\olfileid[hi]{sfr}{siz}{nen}

\olsection{\printtoken{S}{nonenumerable} समुच्चय}

\begin{editorial}
  यह खंड \olref[enm]{sec} की परिभाषा का उपयोग करके $\Bin^\omega$ और
  $\Pow{\PosInt}$ की अगणनीयता सिद्ध करता है। इसे \olref[enm-alt]{sec} के
  रूपांतर से कुछ अधिक प्रारंभिक और कुछ अधिक विस्तृत बनाया गया है।
\end{editorial}

कुछ समुच्चय अनंत होते हैं, जैसे धन पूर्णांकों का समुच्चय $\PosInt$। अब तक
हमने अनंत समुच्चयों के जो उदाहरण देखे हैं, वे सभी !!{enumerable} थे। लेकिन
ऐसे अनंत समुच्चय भी होते हैं जिनमें यह गुण नहीं होता। ऐसे समुच्चय
\emph{!!{nonenumerable}} कहलाते हैं।

सबसे पहले, !!{nonenumerable} समुच्चयों का होना ही शायद आश्चर्यजनक लगे। किसी
भी !!{enumerable} समुच्चय~$A$ के लिए कोई !!a{surjective} फलन
$f \colon \PosInt \to A$ होता है। यदि कोई समुच्चय !!{nonenumerable} है, तो
ऐसा कोई फलन नहीं होता। अर्थात, $\PosInt$ के अनंत !!{element} को~$A$ में भेजने
वाला कोई भी फलन~$A$ के सभी अवयवों तक नहीं पहुँच सकता। इसलिए~$A$ में अनंत धन
पूर्णांकों से भी ``अधिक'' !!{element} होते हैं।

हम यह कैसे सिद्ध करें कि कोई समुच्चय !!{nonenumerable} है? हमें दिखाना होगा
कि ऐसा कोई आच्छादक फलन हो ही नहीं सकता। समान रूप से, हमें दिखाना होगा कि~$A$
के अवयवों को एक दिशा में चलने वाली अनंत सूची में नहीं गिना जा सकता। इसका सबसे
अच्छा ढंग यह दिखाना है कि~$A$ के !!{element} की प्रत्येक सूची में कम-से-कम एक
अवयव अवश्य छूटेगा; अथवा कोई भी फलन $f\colon \PosInt \to A$
!!{surjective} नहीं हो सकता। हम कैंटर की \emph{विकर्ण विधि} से ऐसा कर सकते
हैं।~$A$ के !!{element} की कोई सूची, मान लीजिए $x_1$, $x_2$, \dots, दी हो,
तो हम~$A$ का एक और अवयव इस प्रकार बनाते हैं कि वह अपने निर्माण के कारण उस
सूची में हो ही नहीं सकता।

हमारा पहला उदाहरण समुच्चय~$\Bin^\omega$ है, जिसमें $0$ और $1$ के सभी अनंत,
बिना अंतराल वाले अनुक्रम हैं।

\begin{thm}
\ollabel{thm:nonenum-bin-omega}
$\Bin^\omega$~!!{nonenumerable} है।
\end{thm}

\begin{proof}
विरोधाभास प्राप्त करने के लिए मान लीजिए कि $\Bin^\omega$
!!{enumerable} है; अर्थात मान लीजिए कि सभी !!{element} की कोई सूची
$s_{1}$, $s_{2}$, $s_{3}$, $s_{4}$, \dots{} है और ये~$\Bin^\omega$ के
अवयव हैं। इनमें से प्रत्येक
$s_i$ स्वयं $0$ और~$1$ का एक अनंत अनुक्रम है। इस सूची के किसी अनुक्रम के
$j$वें अवयव को, जब वह सूची में $i$वाँ अनुक्रम हो, $s_i(j)$ कहें। तब $i$वाँ
अनुक्रम~$s_i$ है:
\[
s_i(1), s_i(2), s_i(3), \dots
\]

हम इस सूची और इसमें प्रत्येक अनुक्रम $s_i$ के अवयवों को एक सरणी में व्यवस्थित
कर सकते हैं:
\[
\begin{array}{c|c|c|c|c|c}
& 1 & 2 & 3 & 4 & \dots \\\hline
1 & \mathbf{s_{1}(1)} & s_{1}(2) & s_{1}(3) & s_1(4) & \dots \\\hline
2 & s_{2}(1)& \mathbf{s_{2}(2)} & s_2(3) & s_2(4) & \dots \\\hline
3 & s_{3}(1)& s_{3}(2) & \mathbf{s_3(3)} & s_3(4) & \dots \\\hline
4 & s_{4}(1)& s_{4}(2) & s_4(3) & \mathbf{s_4(4)} & \dots \\\hline
\vdots & \vdots & \vdots & \vdots & \vdots & \mathbf{\ddots}
\end{array}
\]
बगल में नीचे की ओर लिखे क्रमांक सूची के अनुक्रम $s_1$, $s_2$, \dots{} का
क्रम बताते हैं; ऊपर लिखे क्रमांक प्रत्येक अनुक्रम के !!{element} को क्रमांकित
करते हैं। उदाहरण के लिए, $s_{1}(1)$ उस संख्या का नाम है---चाहे वह $0$ हो या
$1$---जो अनुक्रम $s_{1}$ का पहला !!{element} है; और इसी प्रकार आगे।

अब हम एक ऐसा अनंत अनुक्रम $\overline{s}$ बनाएँगे जो $0$ और $1$ से बना है और
इस सूची में हो ही नहीं सकता। $\overline{s}$ की परिभाषा सूची $s_1$, $s_2$, \dots{} पर
निर्भर करेगी। $0$ और $1$ के अनंत अनुक्रमों की कोई भी अनंत सूची एक ऐसा अनंत
अनुक्रम~$\overline{s}$ देती है जिसके सूची में न आने की गारंटी है।

$\overline{s}$ को परिभाषित करने के लिए हम उसके सभी !!{element} बताते हैं;
अर्थात हम $\overline{s}(n)$ बताते हैं, और ऐसा प्रत्येक $n \in \PosInt$ के
लिए करते हैं। हम ऊपर
की सरणी के विकर्ण पर नीचे की ओर पढ़ते हैं (इसीलिए इसका नाम ``विकर्ण विधि''
है), फिर प्रत्येक $1$ को $0$ और प्रत्येक $0$ को~$1$ कर देते हैं। अधिक अमूर्त
रूप में, हम $\overline{s}(n)$ को $0$ या $1$ इस आधार पर परिभाषित करते हैं कि
विकर्ण का $n$वाँ !!{element}, $s_n(n)$, क्रमशः $1$ है या $0$:
\[
\overline{s}(n) =
\begin{cases}
1 & \text{यदि $s_{n}(n) = 0$}\\
0 & \text{यदि $s_{n}(n) = 1$}.
\end{cases}
\]
यदि आपको स्थिति-वार परिभाषा से अधिक सूत्र पसंद हैं, तो आप
$\overline{s}(n) = 1 - s_n(n)$ भी परिभाषित कर सकते हैं।

स्पष्ट है कि $\overline{s}$, $0$ और $1$ का एक अनंत अनुक्रम है, क्योंकि वह
हमारी सरणी के विकर्ण पर आने वाले $0$ और $1$ के अनुक्रम का पूरक अनुक्रम भर
है। अतः $\overline{s}$,~$\Bin^\omega$ का !!a{element} है। लेकिन वह सूची
$s_1$, $s_2$, \dots{} में नहीं हो सकता। क्यों?

वह सूची का पहला अनुक्रम $s_1$ नहीं हो सकता, क्योंकि उसका पहला !!{element}
$s_1$ से भिन्न है। $s_1(1)$ चाहे जो भी हो, हमने $\overline{s}(1)$ को उसका
उलटा परिभाषित किया है। वह सूची का दूसरा अनुक्रम भी नहीं हो सकता, क्योंकि
$\overline{s}$ का दूसरा अवयव $s_2$ के दूसरे अवयव से भिन्न है: यदि $s_2(2)$,
$0$ है, तो $\overline{s}(2)$, $1$ है, और उलटे भी यही होता है। इसी प्रकार आगे।

अधिक ठीक रूप में: यदि $\overline{s}$ सूची में होता, तो कोई $k$ ऐसा होता कि
$\overline{s} = s_{k}$। दो अनुक्रम समान होते हैं यदि और केवल यदि वे प्रत्येक
स्थान पर सहमत हों; अर्थात किसी भी~$n$ के लिए $\overline{s}(n) =
s_{k}(n)$। अतः विशेष रूप से $n = k$ लेने पर $\overline{s}(k) = s_{k}(k)$
होना आवश्यक होता। $s_k(k)$ या तो $0$ है या~$1$। यदि वह $0$ है, तो
$\overline{s}(k)$ को~$1$ होना होगा---हमने $\overline{s}$ को इसी प्रकार
परिभाषित किया है। लेकिन यदि $s_k(k) = 1$ हो, तो $\overline{s}$ की परिभाषा
के कारण फिर $\overline{s}(k) = 0$ होगा। दोनों ही स्थितियों में
$\overline{s}(k) \neq s_{k}(k)$।

हमने यह मानकर आरंभ किया था कि~$\Bin^\omega$ के !!{element} की कोई सूची
$s_1$, $s_2$, \dots{} है। इस सूची से हमने एक अनुक्रम~$\overline{s}$ बनाया
और सिद्ध किया कि वह सूची में नहीं हो सकता। लेकिन वह निश्चित रूप से $0$ और
$1$ का अनुक्रम है, यदि सभी $s_i$, $0$ और $1$ के अनुक्रम हैं; अर्थात
$\overline{s} \in
\Bin^\omega$। इससे विशेष रूप से स्पष्ट होता है कि
~$\Bin^\omega$ के \emph{सभी} !!{element} की कोई सूची नहीं हो सकती, क्योंकि
किसी भी ऐसी सूची से हम एक ऐसा अनुक्रम~$\overline{s}$ बना सकते हैं जिसके सूची
में न होने की गारंटी है। इसलिए~$\Bin^\omega$ के सभी अनुक्रमों की सूची होने
की धारणा विरोधाभास उत्पन्न करती है।
\end{proof}

\begin{explain}
यह प्रमाण-विधि ``विकर्णीकरण'' कहलाती है, क्योंकि यह~$\overline{s}$ को
परिभाषित करने में सरणी के विकर्ण का उपयोग करती है। विकर्णीकरण में सरणी का
होना आवश्यक नहीं है: जब कोई सरणी या वास्तविक विकर्ण न हो, तब भी हम इसी
विचार से दिखा सकते हैं कि समुच्चय !!{enumerable} नहीं हैं।
\end{explain}

\begin{thm}
\ollabel{thm:nonenum-pownat}
$\Pow{\PosInt}$ !!{enumerable} नहीं है।
\end{thm}

\begin{proof}
हम इसी प्रकार आगे बढ़ते हैं: दिखाते हैं कि~$\PosInt$ के उपसमुच्चयों की
प्रत्येक सूची के लिए $\PosInt$ का कोई ऐसा उपसमुच्चय है जो सूची में नहीं हो
सकता। मान लीजिए~$\PosInt$ के उपसमुच्चयों की निम्नलिखित सूची दी गई है:
\[
Z_{1}, Z_{2}, Z_{3}, \dots
\]
अब हम समुच्चय $\overline{Z}$ इस प्रकार परिभाषित करते हैं कि किसी भी
$n \in \PosInt$ के लिए $n \in \overline{Z}$ यदि और केवल यदि
$n \notin Z_{n}$:
\[
\overline{Z} = \Setabs{n \in \PosInt}{n \notin Z_n}
\]
स्पष्ट है कि $\overline{Z}$ धन पूर्णांकों का समुच्चय है, क्योंकि हमारी
धारणा के अनुसार प्रत्येक~$Z_n$ ऐसा ही है; अतः $\overline{Z} \in
\Pow{\PosInt}$। लेकिन $\overline{Z}$ सूची में नहीं हो सकता। यह दिखाने के लिए
हम स्थापित करेंगे कि प्रत्येक $k \in \PosInt$ के लिए $\overline{Z} \neq
Z_k$।

अब कोई भी $k \in \PosInt$ लें। हमने $\overline{Z}$ को इस प्रकार परिभाषित
किया है कि किसी भी $n \in \PosInt$ के लिए $n \in \overline{Z}$ यदि और केवल
यदि $n \notin Z_n$। विशेष रूप से $n=k$ लेने पर $k \in \overline{Z}$ यदि और
केवल यदि $k \notin Z_k$। लेकिन इससे $\overline{Z} \neq Z_k$ सिद्ध होता है,
क्योंकि $k$ एक का !!a{element} है पर दूसरे का नहीं; इसलिए $\overline{Z}$ और
$Z_k$ के !!{element} भिन्न हैं। चूँकि $k$ कोई भी था, $\overline{Z}$ सूची
$Z_1$, $Z_2$, \dots{} में नहीं है।
\end{proof}

\begin{explain}
पिछले प्रमाण में विकर्ण का उल्लेख नहीं था, लेकिन आप निम्न चित्र द्वारा उसमें
विकर्ण देख सकते हैं। कल्पना कीजिए कि समुच्चय $Z_1$, $Z_2$, \dots{} एक सरणी
में लिखे हैं, जहाँ प्रत्येक !!{element}~$j \in Z_i$ को~$j$वें स्तंभ में लिखा
गया है। मान लीजिए उस सूची के पहले चार समुच्चय $\{1,2,3,\dots\}$,
$\{2, 4, 6, \dots\}$, $\{1,2,5\}$ और $\{3,4,5,\dots\}$ हैं। तब सरणी का
आरंभ इस प्रकार होगा:
\[
\begin{array}{r@{}rrrrrrr}
  Z_1 = \{ & \mathbf{1}, & 2, & 3, & 4, & 5, & 6, & \dots\}\\
  Z_2 = \{ &  & \mathbf{2}, &  & 4, &  & 6, & \dots\}\\
  Z_3 = \{ & 1, & 2, &  &  & 5\phantom{,} &  & \}\\
  Z_4 = \{ &  &  & 3, & \mathbf{4}, & 5, & 6, & \dots\}\\
  \vdots & & & & & \ddots
\end{array}
\]
तब विकर्ण पर नीचे की ओर चलते हुए हमें $\overline{Z}$ मिलता है: विकर्ण पर
आने वाली प्रत्येक संख्या को छोड़ देते हैं और उन $j$ को सम्मिलित करते हैं
जहाँ सरणी का $j$वीं पंक्ति/स्तंभ वाला खाना रिक्त है। ऊपर के उदाहरण
में हम $1$ और $2$ को छोड़ेंगे,~$3$ को सम्मिलित करेंगे,~$4$ को छोड़ेंगे, और
इसी प्रकार आगे बढ़ेंगे।
\end{explain}

\begin{prob}
विकर्ण तर्क द्वारा दिखाइए कि $\Pow{\Nat}$ !!{nonenumerable} है।
\end{prob}

\begin{prob}\label{sfr:siz:nen:prob:f-posint}
स्पष्ट विकर्ण तर्क द्वारा दिखाइए कि फलनों $f \colon \PosInt \to \PosInt$
का समुच्चय !!{nonenumerable} है। अर्थात, दिखाइए कि यदि $f_1$, $f_2$, \dots,
फलनों की सूची है और प्रत्येक $f_i\colon
\PosInt \to \PosInt$ है, तो कोई $\overline{f}\colon \PosInt \to
\PosInt$ ऐसा है जो इस सूची में नहीं है।
\end{prob}

\end{document}
